% Lösungen Einheit 2 von 3 – Potenzgesetze (zusammenbau v0.1)
\einheitenkopf{Einheit 2 von 3 · Potenzgesetze}

%% TODO Lösung der Erklärzeile 14g) fehlt in der Bank
\erg{14}{a) gleiche Basis \quad b) $3 \cdot 3 \cdot 3 \cdot 3 \cdot 3$: fünf Faktoren, $3^5 = 243$ \quad c) $4 \cdot 4 \cdot 4 \cdot 4$: vier Faktoren, $4^4 = 256$ \quad d) $5 \cdot 5 \cdot 5 \cdot 5$: vier Faktoren, $5^4 = 625$ \quad e) $2 \cdot 2 \cdot 2 \cdot 2 \cdot 2 \cdot 2 \cdot 2$: sieben Faktoren, $2^7 = 128$ \quad f) $7 \cdot 7 \cdot 7$: drei Faktoren, $7^3 = 343$ \quad g) \ldots \quad h) $5^{4+3} = 5^7 = 78\,125$ \quad i) $7^{8-3} = 7^5 = 16\,807$ \quad j) $2^6 = 64$ \quad k) $(2 \cdot 5)^4 = 10^4 = 10\,000$ \quad l) kein Gesetz: $9 \cdot 8 = 72$ \quad m) $3^4 \cdot 2 = 81 \cdot 2 = 162$ \quad n) $x^7$ \quad o) $10x^7$ \quad p) $3 \cdot 3^4 = 3^5 = 243$ \quad q) $9x^2y^2$ \quad r) $\frac{12x^8}{4x^2} = 3x^6$}
\erg{15}{a) $\frac{1}{25}$}
\erg{16}{a) Die Basis bleibt, nur die Exponenten werden addiert: $3^2 \cdot 3^5 = 3^7 = 2\,187$. \quad b) Als Malkette stehen sechs Dreien hintereinander, also $3^6$; die Basis bleibt, sie wird nicht malgenommen. \quad c) $2^{30} \cdot 2^8 = 2^{38} = 274\,877\,906\,944$ Byte}
